\lesson{3}{Nov 8 2021 Mon (11:28:04)}{The Executive Branch}{Unit 3}

\subsubsection*{The Role of the President and Vice President}

\paragraph{President} The main function of the \bf{Executive Branch} is to carry out and enforce the laws. The \bf{President} is the focus of the \bf{Executive Branch} in the \bf{Constitution}. \it{Article II} describes the qualifications and duties of the \bf{President}. Here are the rules to be \bf{President}:

\begin{itemize}
    \item At least $35$ years old
    \item A U.S. Citizen
    \item A Resident of the \bf{United States} for at least $14$ years.
\end{itemize}

Of course, these are the bare requirements set by the \bf{Constitution}. To win enough votes in the \bf{Electoral College} to become \bf{President}, a candidate must convince voters they are prepared for the difficult duties of the nation's highest office.

\paragraph{Vice President} The \bf{vice President} must meet the same requirements, as the main role of this official is to take over if the \bf{President} is unable to serve. The \bf{vice President} and other \bf{White House staff members} are part of the \bf{Executive Office} of the \bf{President}, or EOP.

Officials in the EOP report directly to the \bf{President} and can serve as advisors on public policy.

\subsubsection*{The Cabinet}

The \bf{Constitution} gives \bf{Congress} the power to create \bf{Executive Departments} and grants the \bf{President} power to appoint officials to lead these departments with approval of the \bf{Senate}. The head of the \bf{Executive Departments} are part of the \bf{President's Cabinet}. They and their staffs manage the daily affairs of policies within their area, such as foreign policy.

By tradition, and not the \bf{Constitution}, the department heads serve as \bf{Advisers} to the \bf{President}. \bf{Cabinet Members} usually have the term \it{Secretary} in their title. Though each member has experience in their field, the \bf{President} has no obligation to follow their guidance. The \bf{President} may choose to allow other officials, such as those within the EOP, to have \it{"Cabinet Status"} meaning they can attend \bf{Cabinet Meetings}, as well as provide counsel to the \bf{President}. Some cabinet members attend meetings of certain EOP offices, as well. The \bf{President's} style of leadership largely determines the relative power of officials in the executive branch. The \bf{President} can remove officials from cabinet leadership positions as desired. \bf{Congress} can create, reorganize, or disband cabinet departments through its policy making process.

\subsubsection*{The Structure of the Executive Branch}

The \bf{Executive Branch} contains many parts. The \bf{President}, \bf{Vice President}, and \bf{Staff in the Executive Office of the President} make up the main part of the \bf{Executive Branch}. Under this main part is the \bf{Cabinet}. The \bf{Cabinet} is a group of leaders of $15$ \bf{Executive Departments} that advise the \bf{President}. Under each \bf{Executive Department} are various \bf{Agencies} with staff that carry out the daily work of enforcing policy.

\bf{Bureaucracy} refers to the:

\begin{itemize}
    \item Officials
    \item Staff
    \item Policies
    \item Regulations
    \item Actions of the Executive Departments and Agencies
\end{itemize}

People differ in how they use the term \bf{Bureaucracy}. Some may use it to refer to all staff working under the \bf{President}, including \bf{Cabinet Members}. Others reserve the term for staff in lower agencies whose jobs mainly are to complete paperwork. In addition, definitions of positions in the \bf{Executive Branch} can be blurry as duties often overlap and depend on how the \bf{President} wishes to organize the work of carrying out policy. About $2,000,000$ people work directly in the \bf{Federal Bureaucracy}. Many more work for it indirectly through partnered independent or state and local agencies and businesses.

\begin{marginfigure}
    \centering
    \includegraphics[width=\textwidth,height=\textheight,keepaspectratio]{Government/images/executive-branch-structure}
    \sidecaption{Executive Branch Structure}
    \label{fig:executive-branch-structure}
\end{marginfigure}

Here's the structure of the \bf{Executive Branch}:

\begin{definition}[President/VP/Executive Office of the President]
    The president heads the executive branch. The EOP includes the vice president and staff in offices such as the OMB (Office of Management and Budget). This office helps the president develop the yearly budget to present to Congress. It oversees executive spending and evaluates the effectiveness of various agencies and programs. Another important part of the EOP is the NSC (National Security Council). Officials with expertise in foreign policy and intelligence gathering, including some cabinet members, meet as part of the NSC. They advise the president on matters of national security and defense.
\end{definition}

\begin{definition}[Department of Agriculture]
    The Department of Agriculture oversees farming activities. It regulates soil conservation efforts and the quality of the food supply. The department publishes nutrition guidance. It also manages programs that assist families such as the school lunch program.
\end{definition}

\begin{definition}[Department of Commerce]
    The Department of Commerce regulates trade, business, and industry. Its mission is to" help make American businesses more innovative at home and more competitive abroad."
\end{definition}

\begin{definition}[Department of Defense]
    The Department of Defense manages all branches of the U.S. military, including the Coast Guard. It contains three million civilian and military employees. The secretary of defense advises the president on weapons systems and troop strength.
\end{definition}

\begin{definition}[Department of Education]
    The Department of Education manages federal financial aid programs. It also collects and distributes research on the nation's education programs. The department enforces equal access to public education.
\end{definition}

\begin{definition}[Department of Energy]
    The Department of Energy regulates the nation's sources of energy, advising the president on many related scientific issues such as developing clean energy. It addresses fuel needs such as oil and nuclear energy.
\end{definition}

\begin{definition}[Department of Health and Human Services]
    The Department of Health and Human Services administers many services Americans use each day. It helps provide for health and other essential needs, especially for those least able to provide for themselves. One in four Americans benefits from Medicare or Medicaid programs. These are health insurance programs provided in close collaboration with state and local governments. The department helps provide important information to Americans on various topics such as diseases and safety. Help for children with disabilities and people suffering from substance abuse are examples of its initiatives. The Food and Drug Administration, or FDA, is an important agency under this department. The FDA regulates food and medicine products, including prescription drugs, vaccines, medical devices, and cosmetics.
\end{definition}

\begin{definition}[Department of Housing and Urban Development]
    The Department Of Housing and Urban Development promotes access to affordable housing for all Americans. Its programs include government mortgages for home ownership, grants to businesses that build and develop communities, enforcement of anti-discrimination laws in housing, and assistance for those facing the loss of their homes.
\end{definition}

\begin{definition}[Department of Homeland Security]
    The Department of Homeland Security regulates immigration policies and efforts to prevent terrorism. It also manages emergency response systems related to natural disaster. This is the youngest of the cabinet departments, created after the events of September 11, 2001. It absorbed the activities of 22 smaller agencies. The Department of Homeland Security is part of the cooperative U.S. Intelligence Community, a group of 16 separate federal agencies responsible for intelligence monitoring, gathering, and reporting to the president and cabinet.
\end{definition}

\begin{definition}[Department of Interior]
    The mission of the Department of Interior is to protect the nation's natural resources and cultural heritage. It regulates the creation and maintenance of federal parks, landmarks, and other federal lands. The department also partners with and provides services to the approximately 595 federally recognized Native American tribes.
\end{definition}

\begin{definition}[Department of Justice]
    The Department of Justice is responsible for prosecuting offenders of federal laws. The head of this department is the Attorney General, who advises the president on preventing crime, punishment for crimes, and ensuring a fair system of justice for all Americans. The Federal Bureau of Investigation, or FBI, is part of the Department of Justice and is responsible for investigating federal crimes and crimes involving more than one state.
\end{definition}

\begin{definition}[Department of Labor]
    The Department of Labor promotes the interests of working Americans, such as the availability of jobs, fair wages, and health and safety while on the job. It also helps people unable to find work and enforces federal laws related to working minors, known as child labor laws, and retirement benefits.
\end{definition}

\begin{definition}[Department of State]
    The Department of State deals with foreign affairs. The Secretary of State is the chief advisor to the president about issues going on in other countries. Diplomacy, the active management of communications with leaders around the world, is considered essential to promoting the interests of Americans, including human rights, security, and prosperity. The U.S. ambassador to the United Nations is part of the Department of State.
\end{definition}

\begin{definition}[Department of Transportation]
    The Department of Transportation (DOT) regulates the safety of, access to, and efficiency of transportation systems such as interstate highways and airplane travel. It monitors and enforces federal laws relating to topics such as the manufacture of vehicles and licensing requirements of operators working for interstate trucking companies.
\end{definition}

\begin{definition}[Department of Treasury]
    The Department of Treasury manages federal finances, including collecting taxes and managing the public debt. It oversees the creation of currency and coins, supervises banking activities, and investigates suspects of tax evasion or counterfeit, which is creating illegal currency.
\end{definition}

\begin{definition}[Department of Veterans Affairs]
    The Department of Veterans Affairs provides many benefits and opportunities to the veterans of the nation's armed forces. Examples include healthcare, education, and home ownership.
\end{definition}

\begin{definition}[Independent Agencies and Government Corporations]
    Congress has also created several dozen independent regulatory agencies and government corporations. These groups exist because they have duties that are beyond the scope of normal policymaking. They are independent in the sense that the president does not have the power to remove its officials at his or her will. Examples include the Federal Reserve, Central Intelligence Agency, and the Federal Communications Commission.
\end{definition}

\subsubsection*{Independent Regulatory Agencies}

\bf{Independent Regulatory Agencies} are offices that serve a specific enforcement function. They operate with much less of the \bf{President's} direction than \bf{Executive Departments}. Like the \bf{Cabinet}, \bf{Congress} creates them and determines through law the powers of the agency. The \bf{President} appoints their \bf{Chief Leaders} with approval from the \bf{Senate}, but usually does not have the power to remove officials without cause.

Reasons for their relative independence include efficiency and democracy. If the functions of these offices depended on legislation or a president's order, their work could be subject to frequent interruption. This could disrupt and undermine time-sensitive services. Another concern is that increasing regulations requires more offices and staff in the executive branch, potentially giving the branch too much power. You may already be familiar with many of these agencies.

\begin{definition}[Federal Communications Commission]
    The Federal Communications Commission regulates communications by television, radio, cable, wire, and satellite.
\end{definition}

\begin{definition}[US Postal Service]
    The US Postal Service regulates and carries out the delivery of the people's mail. It is one of the few agencies mentioned directly in the Constitution.
\end{definition}

\begin{definition}[The Federal Reserve]
    The Federal Reserve is the central bank of the United States. It operates as a partner with the U.S. Treasury to supervise banking services and influence the nation's economy.
\end{definition}

\begin{definition}[The Central Intelligence Agency]
    The Central Intelligence Agency, though it reports to the Director of the U.S. Intelligence community, is more independent than the other agencies that are part of that cooperative. Its main duty is to gather information on foreign governments, businesses, and people. CIA information helps guide the president and Cabinet members on matters of national security.
\end{definition}

\begin{definition}[The Consumer Product Safety Commission]
    The Consumer Product Safety Commission inspects the products we buy for potential safety hazards. When you hear of a recall on a toy, vehicle, or other object, this is usually the work of the CPSC, which has investigated complaints and found the item to be a safety risk.
\end{definition}

\begin{definition}[The Federal Trade Commission]
    The Federal Trade Commission regulates fair business practices. Its goal is to maintain choice for consumers and prevent deception. Individuals can report complaints about businesses to the FTC, such as misleading advertising. The FTC will investigate the complaints and enforce federal law regarding business operations.
\end{definition}

\begin{definition}[The Social Security Administration]
    The Social Security Administration regulates the distribution of social security benefits. Americans who have a disability or are retired as well as family members of a deceased wage earner are eligible for the program. Working Americans pay into the SSA through income taxes.
\end{definition}

\subsubsection*{People Complaining About Big Government}

The executive branch was originally quite small. It has expanded significantly, especially during the Civil War and Great Depression. Congress and the president have increased the size of executive departments and staff to meet expanding needs of the country. The increases in policy mean increases in the need for regulation and enforcement. For all Americans this means more rules to follow in different areas of life, such as owning one's own business, filing income taxes, and receiving assistance after losing a job.

Regulations often mean forms to complete and other documentation to maintain. Many Americans refer to these policies as "red tape" or "jumping through hoops." For many people, the federal bureaucracy has a negative connotation because of the extensive rules to follow as well as its sheer size. Many federal agencies have overlapping areas of supervision, leading to more criticism of government waste and excessive executive power.

However, it is precisely through Americans' demands for services and government accountability that the bureaucracy exists and changes in size. Congress creates legislation to meet the needs and wants of the people. For example, during the Great Depression about one-third of Americans were out of work. They demanded government action to boost the economy and prevent such widespread financial devastation in the future. One of the responses from the federal government was to create the unemployment insurance program. Staff increased in several related federal agencies to administer the program in collaboration with state governments.

\subsubsection*{Congressional Oversight}

\bf{Congress} supervises and can check the \bf{Executive Branch}, including the \bf{Bureaucracy}, in several ways. \bf{Congress} has to approve the existence of \bf{Federal Agencies} through law. The \bf{Executive Branch} cannot spend money on:

\begin{enumerate}
    \item Staffing
    \item Programs
    \item Enforcement
\end{enumerate}

without approval from \bf{Congress} through an \bf{Authorization Bill}, giving permission to start or continue the activities of an agency and setting a maximum amount of money for it. Usually, the actual spending of the money also requires an \bf{Appropriation}, which is permission to spend a specific amount of money for a specific purpose. In most cases, the appropriation is less than the amount authorized. Some programs, such as \bf{Social Security} benefits, operate under a trust fund that is outside the normal federal budget.

The \bf{Legislative Branch} interacts with the \bf{President} and \bf{Federal Agencies} in other ways as well. The \bf{Vice President} is the \bf{President} of the \bf{Senate} when in attendance there. The \bf{President} works with \bf{Congress} to set the policy agenda and try to enact his or her policy goals. Success in achieving executive-designed policy often depends on whether \bf{Congress} is dominated by members of the \bf{President's Political Party}.

\paragraph{Investigating the President}

\bf{Congress} also has the power of investigation over the \bf{Executive Branch}. It is not in the \bf{Constitution} but is considered an \bf{Implied Power}. It might investigate whether complaints that a certain agency's regulations are too costly and ineffective to have merit. \bf{Congress} determines the powers of \bf{Federal Agencies}. It can force \bf{Federal Officials} to appear in court and provide evidence related to activities of different agencies. \bf{Congress} has the power to impeach \bf{Federal Officials}, as well.

Whether \bf{Congress} can force the \bf{President} to submit to a subpoena was the main question under the \bf{Supreme Court} case U.S. v. Nixon. \it{President Richard Nixon} maintained that the \bf{President's} power of \bf{Executive Privilege} allowed him to keep executive records private. His loss in this case caused a release of evidence that led \bf{Congress} to debate whether to impeach him. According to \it{Article II} in the \bf{Constitution}, \bf{Congress} may impeach the \bf{President} for \it{"treason, bribery, or other high crimes or misdemeanors"}. \bf{President Nixon} resigned his position to prevent an impeachment trial, which would have taken a great deal of attention away from important national issues.

\begin{figure}[h]
    \begin{quote}
        \it{"I have never been a quitter. To leave office before my term is completed is abhorrent to every instinct in my body. But as president, I must pull the interests of America first. America needs a full time president and a full time congress, particularly at this time with problems we face at home and abroad. To continue to fight through the months ahead for me personal vindication would almost totally absorb the time and attention of both the president and the congress in a period when our entire focus should be on the great issues of peace abroad and prosperity without inflation at home. Therefore, I shall resign the presidency effective at noon tomorrow. Vice President Ford will be sworn in as president at that hour in this office."}
    \end{quote}
    \caption{President Richard Nixon, Resignation Speech, 1974}
\end{figure}

\subsubsection*{Succession Affecting the Executive Branch}

When a \bf{President's} term of office is up, power must be transferred. The $22$nd Amendment limits the president to two terms of office. This prevents any one person from holding the position indefinitely and increasing its powers. \bf{Presidential} elections happen in November and the swearing in of a new \bf{President} happens in January. The time between is referred to as a \it{"Lame Duck"} period, if the current \bf{President} will be leaving the office. A \it{"Lame Duck President"} may be leaving due to term limits or because he is choosing not to run for reelection or due to losing an election. \it{"Lame Duck Presidents"} often have had limited success in setting policy during those periods.

The transition between two \bf{Presidents} also can lead to a large turnover of staff in the \bf{Executive Branch}. Remember, the \bf{President} has the power to appoint heads of many \bf{Federal Agencies}. The \bf{President} will often reward members of his or her political party for helping the campaign or having exceptional talent in their field. Changes in the administration can lead to great changes in policy and regulation, especially when the new \bf{President} is in a different political party.

First in line for presidential succession is the is the \bf{Vice President}. If the \bf{Vice President} were unable to fulfill the duties of the \bf{President}, next in line is the the \bf{Speaker of the House}. If the \bf{Speaker of the House} is not able to fulfill the duties of \bf{President}, next in line is the the \bf{President Pro Tempore of the Senate}. Next in line are the the \bf{Members of the Cabinet}. The order of succession begins with the \bf{Secretary of State} and moves through all the \bf{Heads of the Executive Departments}.

\newpage
